\chapter[FORMAT DÉFINI PAR L'UTILISATEUR]{Format défini par
l'utilisateur}\label{format-duxe9fini-par-lutilisateur}

APRS ne peut pas anticiper tous les types de données. Le format
User-Defined comble ce vide : les auteurs sont libres d'envoyer
n'importe quel format.

En-tête 3 caractères : - \texttt{\{} : DTI User-Defined. - \texttt{U} :
User ID (1 caractère). - \texttt{X} : type de paquet défini par
l'utilisateur (1 caractère).

\begin{center}
\begin{tabular}{|c|c|c|c|c|}
\hline
\textbf{Champ} & \lit{\{} & User ID & User-Defined Packet Type & User-Defined Data \\
\hline
\textbf{Octets} & 1 & 1 & 1 & n \\
\hline
\end{tabular}
\end{center}

L'APRS Working Group attribue les User IDs sur demande justifiée. Pool
limité.

Pour l'expérimentation ou en attendant l'attribution, tout le monde peut
utiliser \texttt{\{} comme User ID sans notification. Paquets commençant
par \texttt{\{\{} = expérimentaux.

Recommandé : ASCII 7 bits imprimable.

Exemples : - \texttt{\{Q1qwerty} --- User ID \texttt{Q}, type
\texttt{1}. - \texttt{\{\{zasdfg} --- User ID \texttt{\{}
(expérimental), type \texttt{z}.

Ces formats sont tous optionnels. Aucun programme n'est tenu de les
décoder et doit ignorer ceux qu'il ne comprend pas.

\begin{center}\rule{0.5\linewidth}{0.5pt}\end{center}
